% LENS C — Systems / Instrumentation roles
% (Senior Mechatronics & Systems Engineer, precision instrumentation, semiconductor equipment)
\documentclass[10pt,a4paper]{article}
\usepackage[margin=1.6cm]{geometry}
\usepackage{enumitem}
\usepackage{textcomp}
\setlist[itemize]{leftmargin=*, itemsep=1pt, topsep=2pt, parsep=0pt}
\usepackage{titlesec}
\titleformat{\section}{\large\bfseries}{}{0pt}{}[\vspace{-4pt}\rule{\textwidth}{0.4pt}]
\titlespacing*{\section}{0pt}{8pt}{4pt}
\pagestyle{empty}

\newcommand{\Headline}{Senior Mechatronics \& Systems Engineer}

\newcommand{\Summary}{Senior Mechatronics Engineer with 6+ years' experience delivering complex electro-mechanical instruments and systems --- integrating precision mechanics, electronics, motion and control, and software --- from concept to manufacture. I take systems-level ownership: capturing requirements, decomposing them across disciplines, holding the design accountable to specification through structured verification and validation, and managing design change from prototype to production-intent. Experienced in precision mechanism design, electromagnetic/thermal simulation (Motor-CAD), FEA, and cleanroom nanofabrication (ICP-RIE), with direct customer-facing delivery through international on-site commissioning.}

\newcommand{\Competencies}{Systems-level requirements capture \& decomposition \; $\cdot$ \; Precision mechanism \& motion system design \; $\cdot$ \; Verification \& validation against specification \; $\cdot$ \; Tolerance analysis, limits \& fits \; $\cdot$ \; FEA \& electromagnetic/thermal simulation (Motor-CAD) \; $\cdot$ \; Design change control \; $\cdot$ \; Sensors, DAQ \& instrumentation \; $\cdot$ \; Multidiscipline coordination (mechanical, electronics, software, chemistry) \; $\cdot$ \; Engineering documentation \& design reviews \; $\cdot$ \; Customer-facing commissioning}

\newcommand{\AssetCoolBullets}{
  \item Owned the mechanical system design of the Dynamic Coating Application Module (DCAM) --- a precision fluid-delivery and spray-application instrument deployed onto power-grid conductors --- decomposing customer requirements into mechanical, actuation, sensing, and fluid-handling subsystems, from feasibility layouts through detailed design to CNC-machined production.
  \item Held the design accountable to functional requirements and client specification through structured lab test programmes across build iterations (3D-printed concept to production-intent), with tolerance analysis and FMEA across the assembly; managed design change between prototype and production builds.
  \item Designed the fluid-handling path (pumped coating delivery, application heads, flow uniformity), structural mechanics, and actuation to demanding environmental requirements --- outdoor exposure, chemical compatibility, vibration --- coordinating closely with the electronics and software teams on integrating their sensing and control subsystems with the mechanical design.
}

\newcommand{\EducationExtra}{ --- cleanroom nanofabrication project: black silicon via ICP-RIE (SF$_6$/O$_2$)}

\newcommand{\Skills}{%
\textbf{Systems \& Design:} Requirements capture and decomposition, precision mechanisms, motion systems, design change control, work scoping, multidiscipline coordination\\
\textbf{Mechanical Design \& CAD:} SolidWorks (Advanced), Onshape (Advanced), Fusion 360 (Advanced); assembly design, 2D drawing packages to BS~8888 with GD\&T, limits \& fits\\
\textbf{Analysis \& Simulation:} FEA (stress, thermal, dynamic), Motor-CAD (electromagnetic \& thermal), tolerance analysis, MATLAB/Simulink, MuJoCo, OpenSim\\
\textbf{Instrumentation \& Electronics:} Sensors and DAQ, ToF/current/voltage sensing, embedded controllers, BLDC/FOC motor control, CAN protocol work, signal conditioning, PCB collaboration\\
\textbf{Fabrication:} CNC machining, 3D printing, WaterJet; cleanroom nanofabrication (ICP-RIE black silicon, SF$_6$/O$_2$)\\
\textbf{Quality \& Delivery:} FMEA, reliability engineering, verification \& validation, design reviews, BS~8888 and relevant ISO standards; team leadership (up to 5 engineers), mentoring, on-site commissioning}

\newcommand{\Interests}{Precision instrumentation, applied simulation, semiconductor fabrication, verification of learned/physical systems, mentoring engineers.}

\begin{document}
% ============================================================
% SPINE — shared, canonical content. Facts never change here.
% Rulings: decommissioned line (never "live"/"live-line"/"energised");
% canonical titles fixed; only lens files vary emphasis.
% ============================================================

\begin{center}
  {\LARGE\bfseries Uljan Sinani}\\[2pt]
  {\large \Headline{} $|$ United Kingdom}\\[3pt]
  uljan.sinani@gmail.com \; $\cdot$ \; 073-7977-8240 \; $\cdot$ \; linkedin.com/in/uljansinani \; $\cdot$ \; Full UK driving licence\\[2pt]
  {\small Portfolio: \textbf{sinani.ai} — physics-based simulation, world-model latent audit, open engineering notebook \; $\cdot$ \; github.com/USinani}
\end{center}

\section*{Professional Summary}
\Summary

\section*{Core Competencies}
\Competencies

\section*{Professional Experience}

\noindent\textbf{AssetCool Ltd., Leeds, United Kingdom} \hfill May 2025 -- June 2026\\
\textit{Senior Mechatronics Engineer (Mechanical Design Lead)}
\begin{itemize}
  \AssetCoolBullets
  \item Delivered international field commissioning and validation with Hydro-Qu\'ebec, Canada (June 2025) --- full uniform 360\textdegree{} coating across a 300\,m conductor span on a decommissioned line --- including installation health \& safety planning for site work on utility grid infrastructure.
  \item Coordinated a multidiscipline design effort across mechanical, electronics, software, chemistry, and fluidics teams; held the mechanical design accountable to functional requirements and client specification, managing design change between prototype and production-intent builds.
\end{itemize}

\noindent\textbf{Federal-Mogul Powertrain / Tenneco, Essex} \hfill 2023 -- 2025\\
\textit{Mechatronics Design Engineer (Lead)}
\begin{itemize}
  \item Led a five-person team investigating and redesigning the power-electronics water cooling plate on the Belt Integrated Starter-Generator (AMG programme): hydraulic/thermal design, tolerance analysis, and FEA, delivered as a complete 3D model and 2D drawing package to BS~8888 with GD\&T; coordinated the checking, approval, and validation handover (shaker and integration testing) against customer specification.
  \item Owned the mechanical design of the generator housing through a client-driven change programme (new connectors and mounting points for range expansion) --- implementing design change control and reporting progress directly to the end customer alongside internal stakeholders.
  \item Built an electromagnetic/thermal simulation of the switched reluctance motor in Motor-CAD to characterise flux and heat dissipation; used the temperature-gradient insight to define optimal thermistor placement and advise the power-electronics team on power modulation and conductor/harness routing to minimise EMI.
  \item Designed, prototyped, and validated a precision thermistor clamp for repeatable sensor placement within the stator bore, carrying the design from lab verification through to final customer approval at the Munich production site.
  \item Delivered technical drawings to relevant ISO and BS standards for the housing, cooling plate, and rotor; reviewed and approved design work from junior engineers and mentored the team on analysis methods, CAD standards, and documentation practice.
\end{itemize}

\noindent\textbf{Borg Automotive UK Ltd., West Midlands} \hfill 2019 -- 2020\\
\textit{Mechatronics Engineer (R\&D)}
\begin{itemize}
  \item Designed bespoke electro-mechanical test fixtures and automated validation systems for electric power-steering assemblies --- closed-loop test rigs integrating motors, sensors, DAQ, and embedded controllers; conducted stress, thermal, and dynamic FEA to verify sub-assembly robustness under operational loads.
  \item Standardised CAD modelling conventions and 2D drawing standards across the design office; conducted design reviews to improve documentation quality, consistency, and checking efficiency.
\end{itemize}

\noindent\textbf{Car Parts Industries (CPI), Belgium} \hfill 2018 -- 2019\\
\textit{Mechatronics Engineer -- R\&D}
\begin{itemize}
  \item R\&D and New Product Introduction for electric steering columns (VW, FIAT, Opel/Vauxhall); on the Prometheus project, developed an embedded module interfacing with the CAN communication between steering units and the vehicle ECU --- hands-on work with automotive connectors, harnessing, and electrical interconnects.
\end{itemize}

\noindent\textbf{BORG Automotive A/S, Denmark} \hfill 2016 -- 2017\\
\textit{Junior Mechatronics Engineer}
\begin{itemize}
  \item New Product Introduction for electric and mechanical steering racks; authored test-rig specifications, developed jigs and diagnostic test equipment, and fed structured issue reports back to the remanufacturing team.
\end{itemize}

\section*{Doctoral Research}
\noindent\textbf{University of Reading --- PhD, Biomedical Engineering / Collaborative Robotics (distance, part-time)} \hfill Nov 2021 -- Present\\
\textit{Physics-based biomechanical simulation of human upper-limb motion}
\begin{itemize}
  \item Developing two- and three-link rigid-body dynamic models of the human arm in MATLAB with anthropometric parameterisation; integrating Hill-type muscle actuators (OpenSim) and building 3D visualisation pipelines in MuJoCo and Unity --- applying structured modelling, simulation, and performance-evaluation methods to complex mechanical system behaviour.
\end{itemize}

\section*{Education}
\noindent\textbf{Tallinn University of Technology --- MSc, Mechatronics Engineering} \hfill 2014 -- 2016\\
Exchange, University of Southern Denmark --- Mechatronics for Nanotechnology\EducationExtra\\
\textbf{Polytechnic University of Tirana --- BSc, Mechatronics Engineering} \hfill 2012 -- 2014

\section*{Technical Skills}
\Skills

\section*{Courses \& Interests}
Udacity Robotics Software Engineer Nanodegree \; $|$ \; Industrial Automation \; $|$ \; Leadership in Engineering\\
\Interests

\end{document}
